\chapter{Planificación del proyecto}

La planificación del proyecto se organiza entre marzo y diciembre de 2026, contemplando tanto los hitos definidos por la asignatura como las actividades técnicas necesarias para el desarrollo del prototipo. El cronograma articula la elaboración documental, el relevamiento bibliográfico, el user research, la definición de requerimientos, el diseño de arquitectura, la implementación de módulos, la validación y la preparación de la defensa final.

\section{Hitos principales de la asignatura}

\begin{table}[H]
\centering
\caption{Hitos principales de la asignatura y entregables asociados}
\label{tab:hitos_pfi}
\small
\begin{tabularx}{\textwidth}{p{0.24\textwidth} p{0.20\textwidth} X}
\hline
\textbf{Hito} & \textbf{Fecha / período} & \textbf{Detalle} \\
\hline
Inicio de cursada & 10/03/2026 & Presentación de la materia, normativa, instructivo de PFI, propuestas de tema, tutores y áreas temáticas sugeridas. \\
Revisión de propuesta y planificación & 31/03/2026 & Revisión de documento de PFI, explicación de secciones, objetivos, alcance y planificación. \\
Metodología de investigación & 14/04/2026 & Normativa, escritura, gramática, método científico y clase de metodología de investigación. \\
Entrega de propuesta de tema & 25/04/2026 & Fecha límite para la entrega de propuestas de tema. \\
Pautas de presentación preliminar I & 12/05/2026 & Técnicas de research, redacción, notas al pie y bibliografía. \\
Taller de fuentes de investigación & 02/06/2026 & Uso de recursos de biblioteca, fuentes de investigación, bibliografía y herramientas de bibliografía. \\
Entrega avance 25\% & 13/06/2026 & Entrega del avance del 25\%. \\
Oratoria y presentaciones eficaces & 23/06/2026 & Clase orientada a comunicación oral y presentaciones. \\
Feedback y arquitectura & 30/06/2026 & Feedback, pautas de presentación preliminar II y taller de arquitectura. \\
Modelo económico & 07/07/2026 & Taller de modelo económico. \\
Consultas y trabajo en clase & 21/07/2026 & Clase de consultas y trabajo con PFIs en clase. \\
\hline
\end{tabularx}
\end{table}

\begin{table}[H]
\centering
\caption{Hitos principales de la asignatura y entregables asociados (continuación)}
\label{tab:hitos_pfi_continuacion}
\small
\begin{tabularx}{\textwidth}{p{0.24\textwidth} p{0.20\textwidth} X}
\hline
\textbf{Hito} & \textbf{Fecha / período} & \textbf{Detalle} \\
\hline
Exposición avance 50\% & 18/08/2026, 25/08/2026 y 05/09/2026 & Presentación de avance al 50\%, con 10 diapositivas y 10 minutos de exposición oral. La presentación debe subirse hasta el 18/08. \\
Feedback presentación preliminar III & 15/09/2026 & Feedback y pautas para la presentación preliminar III. \\
Modelo de negocio y apoyo visual & 29/09/2026 & Modelo de negocio, herramientas de apoyo visual y consultas. \\
Informe avance 75\% & 20/10/2026 & Entrega del informe escrito con aval del tutor y exposición oral para aprobación de cursada. Fecha recomendada para quienes quieran defender en diciembre. \\
Exposiciones avance 75\% & 27/10/2026 y 03/11/2026 & Exposiciones orales correspondientes al avance del 75\%. El informe escrito y la presentación deben subirse hasta el 20/10. \\
Feedback y consultas & 10/11/2026 & Clase de feedback y consultas. \\
Recuperatorio avance 75\% & 24/11/2026 & Recuperatorio del informe de avance 75\%. \\
Defensas llamado regular & 09/12/2026 al 14/12/2026 & Defensas de PFI en llamado regular. \\
Entrega para defensa en febrero & 09/12/2026 al 16/12/2026 & Período de entrega del PFI para presentación del informe final y defensa en febrero. \\
\hline
\end{tabularx}
\end{table}

\section{Cronograma general de trabajo}

\begin{table}[H]
\centering
\caption{Cronograma general de trabajo del proyecto}
\label{tab:gantt_pfi}
\scriptsize
\setlength{\tabcolsep}{2.2pt}
\begin{tabularx}{\textwidth}{p{0.30\textwidth} c c c c c c c c c c}
\hline
\textbf{Actividad} & \textbf{Mar} & \textbf{Abr} & \textbf{May} & \textbf{Jun} & \textbf{Jul} & \textbf{Ago} & \textbf{Sep} & \textbf{Oct} & \textbf{Nov} & \textbf{Dic} \\
\hline
Definición del tema y propuesta & X & X &  &  &  &  &  &  &  &  \\
Revisión normativa, objetivos y alcance & X & X &  &  &  &  &  &  &  &  \\
Relevamiento bibliográfico inicial &  & X & X & X &  &  &  &  &  &  \\
Estado del arte y marco teórico &  &  & X & X & X &  &  &  &  &  \\
Entrega avance 25\% &  &  &  & X &  &  &  &  &  &  \\
User research exploratorio &  &  &  & X & X & X &  &  &  &  \\
Análisis de feedback profesional &  &  &  & X & X & X &  &  &  &  \\
Definición de requerimientos &  &  &  &  & X & X &  &  &  &  \\
Evaluación de alcance axial complementario &  &  &  &  & X & X & X &  &  &  \\
Diseño de arquitectura del sistema &  &  &  &  & X & X & X &  &  &  \\
Presentación avance 50\% &  &  &  &  &  & X & X &  &  &  \\
Preparación de datasets y preprocesamiento &  &  &  &  &  & X & X &  &  &  \\
Módulo de segmentación sagital &  &  &  &  &  & X & X & X &  &  \\
Mediciones y visualización &  &  &  &  &  &  & X & X &  &  \\
Salida estructurada editable &  &  &  &  &  &  & X & X & X &  \\
Integración del prototipo &  &  &  &  &  &  &  & X & X &  \\
Informe y exposición avance 75\% &  &  &  &  &  &  &  & X & X &  \\
Validación técnica y revisión profesional &  &  &  &  &  &  &  & X & X &  \\
Resultados, discusión y conclusiones &  &  &  &  &  &  &  &  & X & X \\
Corrección final y defensa &  &  &  &  &  &  &  &  &  & X \\
\hline
\end{tabularx}
\end{table}

El cronograma podrá ajustarse en función de los resultados del user research, la disponibilidad de datasets y la factibilidad técnica observada durante la implementación. En particular, la posible incorporación del plano axial se considera una extensión complementaria surgida del relevamiento preliminar con profesionales y del análisis de antecedentes. Por este motivo, su alcance definitivo se consolidará en la entrega del 50\%, manteniendo al plano sagital como núcleo principal del MVP.
